\documentclass[a4paper,12pt]{article}
\usepackage{amsmath,amssymb}
\usepackage{xcolor}
\usepackage[most]{tcolorbox}
\usepackage{geometry}
\geometry{left=2cm, right=2cm, top=2.5cm, bottom=2.5cm}

% Couleurs pour les box
\definecolor{SolutionColor}{HTML}{F0F8FF}
\definecolor{ExerciseColor}{HTML}{E6E6FA}

% Environnements stylisés
\newtcolorbox{solutionbox}[1]{
    colback=SolutionColor,
    colframe=blue!75!black,
    title=#1,
    fonttitle=\bfseries,
    breakable
}

\newtcolorbox{exercice}[1][]{
    colback=ExerciseColor,
    colframe=purple!75!black,
    title=#1,
    fonttitle=\bfseries
}

\begin{document}

% Exercice 7
\begin{exercice}[Exercice 7 $\quad \star \quad$ { [Représenter]}]{}
Pour chaque fonction, tracer la courbe $C_f$ à l'aide de la calculatrice puis déterminer graphiquement l'équation de la tangente à $C_f$ au point d'abscisse $a$.

\begin{enumerate}
    \item $f(x) = -x^3 + 5x + 7$ et $a = 2$,
    \item $f(x) = \dfrac{x}{3 + 8x}$ et $a = 1$,
    \item $f(x) = \sqrt{3x + 7}$ et $a = 3$,
    \item $f(x) = \dfrac{2x + 7}{\sqrt{x}}$ et $a = 2$.
\end{enumerate}
\end{exercice}

\begin{solutionbox}{Solution}
L'idée est de déterminer le nombre dérivé au point $a$ en utilisant le taux d'accroissement, c’est-à-dire la limite du taux :
\[
f'(a) = \lim_{h \to 0} \frac{f(a+h) - f(a)}{h}.
\]

Puis, l'équation de la tangente au point $A$ d'abscisse $a$ est :
\[
y = f(a) + f'(a)(x - a).
\]

\textbf{1. Pour } $f(x) = -x^3 + 5x + 7$ et $a = 2$ :

Calculons le nombre dérivé par le taux d'accroissement :
\[
f(a) = f(2) = -(2)^3 + 5 \cdot 2 + 7 = -8 + 10 + 7 = 9.
\]

\[
f'(2) = \lim_{h \to 0} \frac{f(2+h) - f(2)}{h} = \lim_{h \to 0} \frac{[-(2+h)^3 + 5(2+h) + 7] - 9}{h}.
\]

Développons :
\[
(2+h)^3 = 8 + 12h + 6h^2 + h^3.
\]
\[
f(2+h) = -(8 + 12h + 6h^2 + h^3) + 5(2+h) + 7.
\]
\[
f(2+h) = -8 -12h -6h^2 - h^3 + 10 + 5h + 7 = 9 -7h -6h^2 - h^3.
\]

Donc :
\[
\frac{f(2+h)-f(2)}{h} = \frac{[9 -7h -6h^2 - h^3] - 9}{h} = \frac{-7h -6h^2 - h^3}{h} = -7 - 6h - h^2.
\]

En faisant tendre $h \to 0$ :
\[
f'(2) = -7.
\]

L'équation de la tangente :
\[
y = f(2) + f'(2)(x - 2) = 9 - 7(x - 2) = 9 - 7x + 14 = -7x + 23.
\]

\textbf{Tangente : } $y = -7x + 23$.

\textbf{2. Pour } $f(x) = \dfrac{x}{3 + 8x}$ et $a = 1$ :

\[
f(1) = \frac{1}{3 + 8} = \frac{1}{11}.
\]

Calcul du nombre dérivé :
\[
f'(1) = \lim_{h \to 0} \frac{f(1+h) - f(1)}{h} = \lim_{h \to 0} \frac{\frac{1+h}{3 + 8(1+h)} - \frac{1}{11}}{h}.
\]

Simplifions le dénominateur de $f(1+h)$ :
\[
3+8(1+h)=3+8+8h=11+8h.
\]
Donc :
\[
f(1+h)=\frac{1+h}{11+8h}.
\]

\[
f(1+h)-f(1) = \frac{1+h}{11+8h} - \frac{1}{11}.
\]

Mettons au même dénominateur :
\[
\frac{1+h}{11+8h} - \frac{1}{11} = \frac{11(1+h) - (11+8h)}{11(11+8h)} = \frac{11+11h - 11 -8h}{11(11+8h)} = \frac{3h}{11(11+8h)}.
\]

\[
\frac{f(1+h)-f(1)}{h} = \frac{3h}{h \cdot 11(11+8h)} = \frac{3}{11(11+8h)}.
\]

En faisant $h \to 0$ :
\[
f'(1) = \frac{3}{11 \cdot 11} = \frac{3}{121}.
\]

L'équation de la tangente au point $(1;1/11)$ :
\[
y = \frac{1}{11} + \frac{3}{121}(x - 1).
\]

\textbf{Tangente : } $y = \frac{1}{11} + \frac{3}{121}(x - 1)$.

\textbf{3. Pour } $f(x)=\sqrt{3x+7}$ et $a=3$ :

\[
f(3)=\sqrt{3 \cdot 3+7}=\sqrt{16}=4.
\]

\[
f'(3)=\lim_{h \to 0}\frac{f(3+h)-f(3)}{h}=\lim_{h \to 0}\frac{\sqrt{3(3+h)+7}-4}{h}.
\]

\[
3(3+h)+7=9+3h+7=16+3h.
\]
\[
f(3+h)=\sqrt{16+3h}.
\]

\[
\frac{f(3+h)-f(3)}{h}=\frac{\sqrt{16+3h}-4}{h}.
\]

Pour calculer cette limite, on peut utiliser l'astuce de la conjugaison :
\[
\sqrt{16+3h}-4 = \frac{(\sqrt{16+3h}-4)(\sqrt{16+3h}+4)}{\sqrt{16+3h}+4} = \frac{16+3h-16}{h(\sqrt{16+3h}+4)}=\frac{3h}{h(\sqrt{16+3h}+4)}.
\]

\[
f'(3)=\lim_{h \to 0}\frac{3}{\sqrt{16+3h}+4}=\frac{3}{4+4}= \frac{3}{8}.
\]

Tangente en $(3;4)$ :
\[
y=4+\frac{3}{8}(x-3).
\]

\textbf{Tangente : } $y=4+\frac{3}{8}(x-3)$.

\textbf{4. Pour } $f(x)=\frac{2x+7}{\sqrt{x}}$ et $a=2$ :

\[
f(2)=\frac{2 \cdot 2+7}{\sqrt{2}}=\frac{11}{\sqrt{2}}.
\]

% \[
% f'(2)=\lim_{h \to 0}\frac{f(2+h)-f(2)}{h}.
% \]
% 
% \[
% f(2+h)=\frac{2(2+h)+7}{\sqrt{2+h}}=\frac{4+2h+7}{\sqrt{2+h}}=\frac{11+2h}{\sqrt{2+h}}.
% \]
% 
% \[
% f(2+h)-f(2)=\frac{11+2h}{\sqrt{2+h}} - \frac{11}{\sqrt{2}}.
% \]
% 
% Mettons sous un dénominateur commun $\sqrt{2+h}\sqrt{2}$ :
% \[
% =\frac{(11+2h)\sqrt{2} - 11\sqrt{2+h}}{\sqrt{2+h}\sqrt{2}}.
% \]
% 
% En posant $h=0$, c'est compliqué, utilisons la conjugaison ou le développement en $h$ petit.
% 
% Approche simplifiée : On peut tenter de factoriser par $h$.
% 
% Au voisinage de $h=0$, $\sqrt{2+h}=\sqrt{2}\sqrt{1+\frac{h}{2}} \approx \sqrt{2}(1+\frac{h}{4}-...)$.
% 
% \[
% f(2+h)=\frac{11+2h}{\sqrt{2+h}} \approx \frac{11+2h}{\sqrt{2}(1+\frac{h}{4})}.
% \]
% Approximons en linéaire :
% \[
% \frac{1}{1+\frac{h}{4}} \approx 1 - \frac{h}{4}.
% \]
% 
% \[
% f(2+h)\approx\frac{11+2h}{\sqrt{2}}(1 - \frac{h}{4})=\frac{11+2h}{\sqrt{2}} - \frac{11+2h}{4\sqrt{2}}h.
% \]
% 
% C'est un peu lourd sans la dérivée connue. Alternativement, essayons de réécrire $f(x)$ :

\[
f(x)=\frac{2x+7}{\sqrt{x}}=\frac{2x}{\sqrt{x}}+\frac{7}{\sqrt{x}}=2\sqrt{x}+\frac{7}{\sqrt{x}}.
\]

Calculons le taux d'accroissement entre $2$ et $2+h$ :

\[
f(2+h)-f(2)=2\sqrt{2+h}+ \frac{7}{\sqrt{2+h}} - \left(2\sqrt{2}+\frac{7}{\sqrt{2}}\right).
\]

Utilisons la dérivée connue des fonctions racine et inverse, mais en mode taux :

Pour gagner du temps :
- $\sqrt{x}$ dérivée en $2$ est $\frac{1}{2\sqrt{2}}$.
- $\frac{7}{\sqrt{x}}=7x^{-1/2}$ dérivée en $2$ est $7(-\frac{1}{2})2^{-3/2}= -\frac{7}{2 \cdot 2^{3/2}}=-\frac{7}{4\sqrt{2}}$.

Somme des dérivées en $2$ :
\[
f'(2)=2 \cdot \frac{1}{2\sqrt{2}} + (-\frac{7}{4\sqrt{2}}) = \frac{1}{\sqrt{2}} - \frac{7}{4\sqrt{2}}=\frac{4-7}{4\sqrt{2}}=\frac{-3}{4\sqrt{2}}.
\]

Tangente en $(2; \frac{11}{\sqrt{2}})$ :
\[
y=\frac{11}{\sqrt{2}}+\frac{-3}{4\sqrt{2}}(x-2).
\]

\textbf{Tangente :} $y=\frac{11}{\sqrt{2}} - \frac{3}{4\sqrt{2}}(x-2)$.

\end{solutionbox}

% Exercice 8
\begin{exercice}[title=Exercice 8 $\quad \star \quad$ { [Calculer]}]
Soit $f$ une fonction définie et dérivable sur $\mathbb{R}$ telle que $f(2)=5$ et le nombre dérivé $f'(2)=-3$. Déterminer l'équation de la tangente à la courbe $C_f$ au point d'abscisse $2$.
\end{exercice}

\begin{solutionbox}{Solution}
Le nombre dérivé $f'(2)=-3$ est déjà donné. Au point $a=2$, $f(2)=5$.

L’équation de la tangente :
\[
y=f(2)+f'(2)(x-2)=5-3(x-2)=5-3x+6=11-3x.
\]

\textbf{Tangente : } $y=-3x+11$.

\end{solutionbox}

% Exercice 9
\begin{exercice}[title=Exercice 9 $\quad \star \quad$ { [Calculer]}]
Soit $f$ définie sur $\mathbb{R}_+^*$ par $f(x)=\sqrt{x}$. On admet que $f$ est dérivable en $2$ et que le nombre dérivé $f'(2)=\frac{1}{4}$. Déterminer l’équation de la tangente à $C_f$ au point d’abscisse $2$.
\end{exercice}

\begin{solutionbox}{Solution}
On sait $f(2)=\sqrt{2}$ et $f'(2)=\frac{1}{4}$.

Tangente :
\[
y = f(2)+f'(2)(x-2)=\sqrt{2}+\frac{1}{4}(x-2).
\]

\textbf{Tangente : } $y=\sqrt{2}+\frac{1}{4}(x-2)$.

\end{solutionbox}

% Exercice 10
\begin{exercice}[title=Exercice 10 $\quad \star \star \quad$ { [Calculer, Représenter]}]
Soit $f$ définie sur $\mathbb{R}_+^*$ par $f(x)=\frac{1}{x}$.
\begin{enumerate}
    \item En calculant le taux de variation, montrer que $f$ est dérivable en $1$ et calculer $f'(1)$.
    \item En déduire l'équation de la tangente $T_1$ à $C_f$ au point d'abscisse $1$.
    \item Étudier la position relative de $T_1$ par rapport à $C_f$.
    \item En procédant comme précédemment, déterminer l'équation de la tangente $T_{-1}$ à $C_f$ au point d’abscisse $-1$. Que peut-on dire des droites $T_1$ et $T_{-1}$ ?
\end{enumerate}
\end{exercice}

\begin{solutionbox}{Solution}
\textbf{1. Calcul du nombre dérivé en 1 :}

\[
f'(1)=\lim_{h \to 0}\frac{f(1+h)-f(1)}{h}=\lim_{h \to 0}\frac{\frac{1}{1+h}-1}{h}.
\]

\[
f(1)=1.
\]

\[
\frac{1}{1+h}-1=\frac{1-(1+h)}{1+h}=\frac{-h}{1+h}.
\]

\[
\frac{f(1+h)-f(1)}{h}=\frac{-h/(1+h)}{h}=\frac{-1}{1+h}.
\]

En faisant $h \to 0$ :
\[
f'(1)=-1.
\]

\textbf{2. Tangente en $x=1$:}

$f(1)=1$ et $f'(1)=-1$.

Tangente :
\[
y=1-1(x-1)=1-(x-1)=2 - x.
\]

\textbf{Tangente $T_1$: } $y=-x+2$.

\textbf{3. Position de $T_1$ par rapport à $C_f$:}

$f(x)=\frac{1}{x}$ et $T_1(y)=2 - x$.

Comparons $f(x)$ et $2 - x$ pour $x>0$ :
\[
\frac{1}{x} \quad \text{et} \quad 2-x.
\]

Pour $x>0$, si $x<1$, $2-x>1$, mais $1/x>1$ aussi. Vérifions en détail :

À $x=1$, $f(1)=1$ et $T_1(1)=1$ se touchent.

Pour $x<1$, $2-x>1$, et $1/x>1$ mais $1/x$ grandit en descendant $x$, on peut tester $x=1/2$ :
$f(1/2)=2$, $T_1(1/2)=2-0.5=1.5$, donc $f(x)>T_1(x)$ pour $x=1/2$.

Pour $x>1$, par exemple $x=2$: 
$f(2)=0.5$, $T_1(2)=0$, donc $f(x)>T_1(x)$ ? Non, ici $f(2)=0.5$ et $T_1(2)=0$ c'est $f(x)>T_1(x)$.

En fait, pour $x>1$, $2-x<1$, mais $1/x<1$, il faut voir lequel est plus grand. À $x=2$, $1/2=0.5$ et $2-2=0$, $f(2)>T_1(2)$.

On constate que pour $x>1$, $f(x)$ descend plus vite. En réalité, la tangente est une corde "approchante" : Comme $f'(1)=-1$, la tangente est sous la courbe si la courbe descend moins vite.

Test $x=0.5$: $f(0.5)=2$, $T_1(0.5)=1.5$, $f>T_1$.  
$x=2$: $f(2)=0.5$, $T_1(2)=0$, $f>T_1$.

Partout, $f(x)>T_1(x)$ sauf en $x=1$ où ils sont égaux.

\textbf{4. Tangente en $x=-1$:}

On refait le même processus :

\[
f'( -1)=\lim_{h \to 0}\frac{f(-1+h)-f(-1)}{h}=\lim_{h \to 0}\frac{\frac{1}{-1+h}-\frac{1}{-1}}{h}.
\]

\[
f(-1)=-1, \quad \frac{1}{-1}= -1.
\]

\[
\frac{1}{-1+h}-\frac{1}{-1}=\frac{-1}{-1+h} - (-1)=\frac{-1}{-1+h}+1=\frac{-1+(-1+h)}{-1+h}=\frac{h}{-1+h}.
\]

\[
\frac{f(-1+h)-f(-1)}{h}=\frac{h/(-1+h)}{h}=\frac{1}{-1+h}.
\]

En faisant $h \to 0$ :
\[
f'(-1)=\frac{1}{-1+0}=\frac{1}{-1}=-1.
\]

Tangente en $x=-1$:
$f(-1)=f(-1)=\frac{1}{-1}=-1$.
$f'(-1)=-1$.

Tangente :
\[
y=f(-1)+f'(-1)(x+1)=-1-1(x+1)=-1-(x+1)=-x-2.
\]

\textbf{Tangente $T_{-1}$ :} $y=-x-2$.

Remarquons que $T_1$ : $y=-x+2$ et $T_{-1}$ : $y=-x-2$. Ces deux droites sont parallèles (même coefficient directeur) et distantes l'une de l'autre.

\end{solutionbox}

% Exercice 16
\begin{exercice}[title=Exercice 16 $\quad \star \star \quad$ { [Représenter, Calculer]}]
Soit $f(x)=ax^2+bx+c$ un polynôme du second degré. On note $C_f$ sa courbe et on considère deux points $A$ et $B$ de $C_f$ d’abscisses $\alpha$ et $\beta$.
\begin{enumerate}
    \item[(a)] Démontrer que l’équation de la tangente à $C_f$ au point $A$ est :
    \[
    y = (2a\alpha + b)x + c - a\alpha^2.
    \]
    \item[(b)] Déterminer l’équation de la tangente à $C_f$ au point $B$.
    \item Démontrer que ces tangentes se coupent en un point d’abscisse $\frac{\alpha + \beta}{2}$.
    \item Vérifier que :
    \[
    f(\beta) - f(\alpha) = (\beta - \alpha) f'\left(\frac{\alpha + \beta}{2}\right).
    \]
    \item En déduire que la tangente à $C_f$ au point d’abscisse $\dfrac{\alpha + \beta}{2}$ est parallèle à la droite $(AB)$.
    \item Faire une figure illustrant le résultat.
\end{enumerate}
\end{exercice}

\begin{solutionbox}{Solution}
\textbf{Idée générale :} Le nombre dérivé de $f$ en $x=\alpha$ se calcule par le taux d'accroissement:
\[
f'(\alpha)=\lim_{h \to 0}\frac{f(\alpha+h)-f(\alpha)}{h}.
\]

Pour un polynôme $f(x)=ax^2+bx+c$ :
\[
f(\alpha)=a\alpha^2+b\alpha+c.
\]

\[
f(\alpha+h)=a(\alpha+h)^2+b(\alpha+h)+c = a(\alpha^2+2\alpha h +h^2)+b(\alpha + h)+c.
\]

\[
f(\alpha+h)-f(\alpha)= a(\alpha^2+2\alpha h +h^2)+b(\alpha + h)+c - (a\alpha^2+b\alpha+c).
\]

Simplifions :
\[
= a(2\alpha h + h^2)+ b h = 2a\alpha h + ah^2 + bh.
\]

\[
\frac{f(\alpha+h)-f(\alpha)}{h} = 2a\alpha + ah + b.
\]

En faisant $h \to 0$ :
\[
f'(\alpha)=2a\alpha + b.
\]

\textbf{(a)} Équation de la tangente en $A(\alpha; f(\alpha))$ :
Le coefficient directeur est $f'(\alpha)=2a\alpha + b$.
\[
y=f(\alpha)+f'(\alpha)(x-\alpha)=(a\alpha^2+b\alpha+c)+(2a\alpha+b)(x-\alpha).
\]

Développons :
\[
y=a\alpha^2+b\alpha+c+(2a\alpha+b)x - (2a\alpha+b)\alpha
\]

\[
y=a\alpha^2+b\alpha+c+(2a\alpha+b)x - (2a\alpha^2+b\alpha)
\]

Simplifions :
\[
y=(2a\alpha+b)x + [a\alpha^2+b\alpha+c -2a\alpha^2 - b\alpha]
\]

\[
y=(2a\alpha+b)x + (c - a\alpha^2)
\]

\textbf{(b)} De même, la tangente en $B(\beta;f(\beta))$ a pour équation :
\[
y=(2a\beta+b)x + c - a\beta^2.
\]

\textbf{(c)} Pour trouver le point d'intersection des deux tangentes, résolvons:
\[
(2a\alpha+b)x + c - a\alpha^2 = (2a\beta+b)x + c - a\beta^2.
\]

\[
(2a\alpha+b)x - (2a\beta+b)x = -a\beta^2 + a\alpha^2
\]

\[
(2a\alpha - 2a\beta)x + (b-b)x = a(\alpha^2 - \beta^2)
\]

\[
2a(\alpha-\beta)x = a(\alpha+\beta)(\alpha-\beta)
\]

\[
2a(\alpha-\beta)x = a(\alpha-\beta)(\alpha+\beta)
\]

Si $\alpha \neq \beta$, on simplifie par $(\alpha-\beta)$:
\[
2a x = a(\alpha+\beta)
\]

\[
x=\frac{\alpha+\beta}{2}.
\]

\textbf{(d)} Vérifions la relation:
\[
f(\beta)-f(\alpha)=(\beta-\alpha)f'\left(\frac{\alpha+\beta}{2}\right).
\]

Calculons $f(\beta)-f(\alpha)$:
\[
f(\beta)=a\beta^2+b\beta+c,\quad f(\alpha)=a\alpha^2+b\alpha+c.
\]

\[
f(\beta)-f(\alpha)=a(\beta^2-\alpha^2)+b(\beta-\alpha) = (\beta-\alpha)(a(\beta+\alpha)+b).
\]

Or, $f'\left(\frac{\alpha+\beta}{2}\right)=2a\cdot\frac{\alpha+\beta}{2}+b=a(\alpha+\beta)+b$.

Donc :
\[
f(\beta)-f(\alpha)=(\beta-\alpha)(a(\alpha+\beta)+b) = (\beta-\alpha)f'\left(\frac{\alpha+\beta}{2}\right).
\]

\textbf{(e)} Interprétons ce résultat: $(AB)$ est la sécante joignant $(\alpha,f(\alpha))$ et $(\beta,f(\beta))$. Le résultat montre que le coefficient directeur de cette sécante est $f'\left(\frac{\alpha+\beta}{2}\right)$. Donc la tangente en $\frac{\alpha+\beta}{2}$ a le même coefficient directeur et est donc parallèle à $(AB)$.

\textbf{(f)} Faire une figure : On dessinera une parabole, deux points $A$ et $B$ dessus, la sécante $(AB)$, et la tangente au point d'abscisse $\frac{\alpha+\beta}{2}$ montrant qu'elle est parallèle à $(AB)$.

\end{solutionbox}

\end{document}
